\documentclass[11pt]{article}

% ---------------------------------------------------------------------------
% Packages (kept deliberately standard so the document compiles with a vanilla
% TeX Live / MiKTeX installation -- pdflatex, two passes for the ToC/refs).
% ---------------------------------------------------------------------------
\usepackage[margin=1in]{geometry}
\usepackage{amsmath,amssymb,amsthm}
\usepackage{bm}
\usepackage{booktabs}
\usepackage{enumitem}
\usepackage[hidelinks]{hyperref}

% ---------------------------------------------------------------------------
% Notation macros
% ---------------------------------------------------------------------------
\newcommand{\R}{\mathbb{R}}
\newcommand{\E}{\mathbb{E}}
\newcommand{\Var}{\operatorname{Var}}
\newcommand{\Cov}{\operatorname{Cov}}
\newcommand{\Normal}{\mathcal{N}}
\newcommand{\KL}{\mathrm{KL}}
\newcommand{\Tr}{^{\mathsf{T}}}              % transpose
\newcommand{\ind}{\mathbf{1}}                % indicator
\newcommand{\given}{\,\vert\,}
\newcommand{\eqdef}{:=}
\DeclareMathOperator*{\argmax}{arg\,max}
\DeclareMathOperator*{\argmin}{arg\,min}
\newcommand{\state}{S}
\newcommand{\obs}{r}
\newcommand{\nstate}{K}
\newcommand{\fwd}{\alpha}
\newcommand{\bwd}{\beta}

\theoremstyle{plain}
\newtheorem{proposition}{Proposition}
\newtheorem{lemma}{Lemma}
\theoremstyle{definition}
\newtheorem{definition}{Definition}
\theoremstyle{remark}
\newtheorem{remark}{Remark}

\title{\textbf{A Gaussian Hidden Markov Model for Market-Regime\\ Detection and Regime-Aware Position Sizing}\\[4pt]
\large Mathematical specification of the \texttt{lng\_desk.hmm} module}
\author{LNG Desk --- Quantitative Research}
\date{\today}

\begin{document}
\maketitle

\begin{abstract}
We give a complete, self-contained mathematical description of a three-state
Gaussian hidden Markov model (HMM) used to classify daily asset returns into
latent \emph{bear}, \emph{neutral} and \emph{bull} regimes, and to convert the
inferred regime belief into a trading position. The document is written for a
reader at the level of a final-year undergraduate in mathematics: we assume
familiarity with probability, linear algebra, and multivariable calculus, but we
derive every algorithm --- the forward--backward recursions, the Baum--Welch
(EM) parameter estimator, the Viterbi decoder, online filtering, and the
predictive moments --- from first principles. Each modelling choice in the
implementation is stated explicitly and justified. The notation maps one-to-one
onto the code in \texttt{src/lng\_desk/hmm/}; a symbol dictionary is given in
Section~\ref{sec:symbols}.
\end{abstract}

\tableofcontents

% ===========================================================================
\section{Motivation and overview}
% ===========================================================================

Daily returns of a traded asset are notoriously hard to model with a single
fixed distribution. Empirically, financial markets alternate between
qualitatively different \emph{regimes}: calm trending periods with small
positive drift and low volatility; directionless choppy periods; and stressed
periods with negative drift and large volatility. These regimes are not labelled
in the data --- we observe only the returns. This is precisely the setting of a
\emph{hidden Markov model}: an unobserved (``hidden'') discrete state evolves as
a Markov chain, and at each time the state emits an observation through a
state-dependent probability law.

We posit $\nstate=3$ hidden states with the interpretation
\[
\text{state }1=\text{bear},\qquad
\text{state }2=\text{neutral},\qquad
\text{state }3=\text{bull},
\]
and we assume that, \emph{conditional on the current state}, the daily return is
Gaussian with a state-specific mean and variance. The bear state carries a
negative mean and (typically) the largest variance; the neutral state a
near-zero mean and small variance; the bull state a positive mean.

The work splits into four classical problems, which we treat in turn:
\begin{enumerate}[label=(\Roman*),leftmargin=2.2em]
  \item \textbf{Evaluation.} Given parameters, compute the likelihood of an
        observed return series. Solved by the \emph{forward} recursion
        (Section~\ref{sec:fb}).
  \item \textbf{Learning.} Estimate the parameters from data by maximum
        likelihood. Solved by the \emph{Baum--Welch} algorithm, an instance of
        Expectation--Maximisation (Section~\ref{sec:em}).
  \item \textbf{Decoding.} Recover the most likely sequence of hidden states.
        Solved by the \emph{Viterbi} algorithm (Section~\ref{sec:viterbi}).
  \item \textbf{Prediction \& control.} Online, using only past data, compute
        the current regime belief, predict tomorrow's return distribution, and
        map it to a position (Sections~\ref{sec:filtering}--\ref{sec:trading}).
\end{enumerate}

% ===========================================================================
\section{The probabilistic model}\label{sec:model}
% ===========================================================================

\subsection{Random variables and indices}

Time is discrete, $t=1,2,\dots,T$, indexing trading days.\footnote{The
implementation uses $0$-based indexing ($t=0,\dots,T-1$); this is immaterial to
the mathematics.} At each $t$ there is
\begin{itemize}[leftmargin=1.6em]
  \item a \textbf{hidden state} $\state_t \in \{1,\dots,\nstate\}$, with
        $\nstate=3$; and
  \item an \textbf{observation} $\obs_t \in \R$, the day's return (e.g.\ the log
        return $\obs_t=\log P_t-\log P_{t-1}$ of the underlying).
\end{itemize}
We write $\obs_{1:t} \eqdef (\obs_1,\dots,\obs_t)$ for a prefix of the
observation sequence, and similarly $\state_{1:t}$.

\subsection{Defining assumptions}

The model is characterised by two conditional-independence assumptions.

\begin{definition}[Markov property of the hidden chain]
The state sequence $(\state_t)$ is a first-order, time-homogeneous Markov chain:
\begin{equation}
  P(\state_t \given \state_{1:t-1}) = P(\state_t\given \state_{t-1}),
  \qquad t\ge 2,
\end{equation}
and this conditional law does not depend on $t$.
\end{definition}

\begin{definition}[Conditional independence of emissions]
Given the entire state path, the observations are mutually independent and each
$\obs_t$ depends only on the contemporaneous state $\state_t$:
\begin{equation}
  P(\obs_{1:T}\given \state_{1:T})
  = \prod_{t=1}^{T} P(\obs_t \given \state_t).
\end{equation}
\end{definition}

\subsection{Parameters}

The model is fully specified by the triple $\theta=(\bm\pi, A, \{(\mu_k,\sigma_k^2)\})$:
\begin{align}
  \pi_i        &\eqdef P(\state_1=i),
                 &&\text{initial distribution, } \bm\pi\in\R^{\nstate},\ \textstyle\sum_i\pi_i=1;\\
  A_{ij}       &\eqdef P(\state_{t+1}=j\given \state_t=i),
                 &&\text{transition matrix, } A\in\R^{\nstate\times\nstate},\ \textstyle\sum_j A_{ij}=1;\\
  b_k(\obs)    &\eqdef \Normal(\obs;\mu_k,\sigma_k^2)
                 = \frac{1}{\sqrt{2\pi\sigma_k^2}}
                   \exp\!\Big(-\tfrac{(\obs-\mu_k)^2}{2\sigma_k^2}\Big),
                 &&\text{emission density of state } k.
\end{align}
Here $A$ is \emph{row-stochastic}: row $i$ is the probability vector of the next
state given current state $i$. We collect the emission parameters as
$\mu=(\mu_1,\dots,\mu_\nstate)\in\R^{\nstate}$ and
$\sigma^2=(\sigma_1^2,\dots,\sigma_\nstate^2)\in\R_{>0}^{\nstate}$.

\subsection{The joint law and the likelihood}

Combining the two assumptions, the joint density of a state path and the
observations factorises along the chain:
\begin{equation}\label{eq:joint}
  P(\state_{1:T},\obs_{1:T}\given\theta)
  = \pi_{\state_1}\, b_{\state_1}(\obs_1)
    \prod_{t=2}^{T} A_{\state_{t-1}\state_t}\, b_{\state_t}(\obs_t).
\end{equation}
The quantity we ultimately want to maximise over $\theta$ is the
\emph{(marginal) likelihood} of the observed data, obtained by summing
\eqref{eq:joint} over all $\nstate^{T}$ possible state paths:
\begin{equation}\label{eq:like}
  L(\theta) \eqdef P(\obs_{1:T}\given\theta)
  = \sum_{\state_{1:T}\in\{1,\dots,\nstate\}^T}
    \pi_{\state_1} b_{\state_1}(\obs_1)
    \prod_{t=2}^{T} A_{\state_{t-1}\state_t} b_{\state_t}(\obs_t).
\end{equation}

A direct evaluation of \eqref{eq:like} is hopeless: the number of paths grows as
$\nstate^{T}$ (for one year of daily data, $3^{252}\approx 10^{120}$). The key
structural insight --- exploited by every algorithm below --- is that the
sum can be reorganised into $T$ steps of $O(\nstate^2)$ work each by
\emph{dynamic programming}, because the chain only ever ``remembers'' the most
recent state.

% ===========================================================================
\section{Evaluation: the forward--backward algorithm}\label{sec:fb}
% ===========================================================================

\subsection{Forward variables}

\begin{definition}[Forward variable]
For each $t$ and state $i$ define the joint density of the data seen so far and
the current state,
\begin{equation}
  \fwd_t(i) \eqdef P(\obs_{1:t},\, \state_t=i \given \theta).
\end{equation}
\end{definition}

\begin{proposition}[Forward recursion]\label{prop:fwd}
The forward variables satisfy
\begin{align}
  \fwd_1(i) &= \pi_i\, b_i(\obs_1), \label{eq:fwd-init}\\
  \fwd_t(j) &= \Big(\sum_{i=1}^{\nstate} \fwd_{t-1}(i)\, A_{ij}\Big) b_j(\obs_t),
              \qquad t=2,\dots,T,\label{eq:fwd-rec}
\end{align}
and the likelihood is recovered by marginalising the last step,
\begin{equation}\label{eq:like-from-fwd}
  L(\theta) = \sum_{i=1}^{\nstate} \fwd_T(i).
\end{equation}
\end{proposition}

\begin{proof}
The initialisation \eqref{eq:fwd-init} is the definition at $t=1$ using
$P(\obs_1,\state_1=i)=P(\state_1=i)P(\obs_1\given\state_1=i)$. For the recursion,
condition on the previous state and use the law of total probability:
\begin{align*}
  \fwd_t(j)
  &= P(\obs_{1:t}, \state_t=j)
   = \sum_{i} P(\obs_{1:t}, \state_{t-1}=i, \state_t=j)\\
  &= \sum_{i} P(\obs_{1:t-1},\state_{t-1}=i)\,
             P(\state_t=j\given \state_{t-1}=i)\,
             P(\obs_t\given \state_t=j),
\end{align*}
where the factorisation uses the Markov property (the transition to $j$ depends
only on $i$) and conditional independence of emissions ($\obs_t$ depends only on
$\state_t=j$, and is independent of $\obs_{1:t-1}$ given $\state_t$). Recognising
the three factors as $\fwd_{t-1}(i)$, $A_{ij}$ and $b_j(\obs_t)$ gives
\eqref{eq:fwd-rec}. Finally \eqref{eq:like-from-fwd} is
$\sum_i P(\obs_{1:T},\state_T=i)=P(\obs_{1:T})$.
\end{proof}

Each step \eqref{eq:fwd-rec} costs $O(\nstate^2)$, so the whole forward pass is
$O(\nstate^2 T)$ --- linear in the length of the series. This single fact is what
makes HMMs tractable.

\subsection{Backward variables}

To learn parameters we also need the \emph{posterior} distribution of the states
given \emph{all} the data. This requires a complementary recursion that
propagates information backwards in time. The recursions below all rest on a
single conditional-independence fact, which we isolate once.

\begin{lemma}[Markov blanket]\label{lem:blanket}
Under the two model assumptions of Section~\ref{sec:model}, the future is
independent of the past given the present state: for every $t$,
\begin{equation}\label{eq:blanket}
  \big(\obs_{t+1:T},\state_{t+1:T}\big) \ \perp\ \big(\obs_{1:t-1},\state_{1:t-1}\big)
  \quad\text{given } \state_t .
\end{equation}
In particular $P(\obs_{t+1:T}\given\state_t,\obs_{1:t})=P(\obs_{t+1:T}\given\state_t)$.
\end{lemma}

\begin{proof}
Fix $\state_t=i$. The joint factorisation \eqref{eq:joint} splits as a product
of a ``past'' block involving only $(\state_{1:t},\obs_{1:t})$ and a ``future''
block $\prod_{u>t} A_{\state_{u-1}\state_u} b_{\state_u}(\obs_u)$ that, given
$\state_t$, involves only indices $u>t$. Two functions of disjoint coordinate
blocks that share only the conditioning variable $\state_t$ are conditionally
independent given $\state_t$; summing the past block over $\state_{1:t-1}$ and
the future block over $\state_{t+1:T}$ leaves \eqref{eq:blanket}. The displayed
consequence is the special case obtained by marginalising the future states.
\end{proof}

\begin{definition}[Backward variable]
\begin{equation}
  \bwd_t(i) \eqdef P(\obs_{t+1:T}\given \state_t=i,\theta),
\end{equation}
the likelihood of the \emph{future} observations given the present state, with
the convention $\bwd_T(i)=1$ for all $i$ (there is no future to condition).
\end{definition}

\begin{proposition}[Backward recursion]\label{prop:bwd}
\begin{align}
  \bwd_T(i) &= 1, \\
  \bwd_t(i) &= \sum_{j=1}^{\nstate} A_{ij}\, b_j(\obs_{t+1})\, \bwd_{t+1}(j),
              \qquad t=T-1,\dots,1.\label{eq:bwd-rec}
\end{align}
\end{proposition}

\begin{proof}
Condition on the next state $j$:
\[
  \bwd_t(i) = \sum_j P(\state_{t+1}=j,\obs_{t+1},\obs_{t+2:T}\given\state_t=i).
\]
Factor the summand using the Markov property and conditional independence:
$P(\state_{t+1}=j\given\state_t=i)=A_{ij}$, then
$P(\obs_{t+1}\given\state_{t+1}=j)=b_j(\obs_{t+1})$, and finally
$P(\obs_{t+2:T}\given\state_{t+1}=j,\state_t=i,\obs_{t+1})
=P(\obs_{t+2:T}\given\state_{t+1}=j)=\bwd_{t+1}(j)$ by the Markov-blanket
Lemma~\ref{lem:blanket} (the future beyond $t+1$ is independent of $\state_t$ and
$\obs_{t+1}$ given $\state_{t+1}$). This is \eqref{eq:bwd-rec}.
\end{proof}

\subsection{Posterior marginals and pairwise posteriors}

The forward and backward variables combine to give the two posterior quantities
that drive learning. Throughout, $L=L(\theta)=\sum_i \fwd_T(i)$.

\begin{proposition}[Smoothed state posterior]\label{prop:gamma}
The posterior probability of being in state $i$ at time $t$ given all data is
\begin{equation}\label{eq:gamma}
  \gamma_t(i)\eqdef P(\state_t=i\given\obs_{1:T},\theta)
  = \frac{\fwd_t(i)\,\bwd_t(i)}{L}.
\end{equation}
\end{proposition}

\begin{proof}
By definition of the forward and backward variables and the Markov-blanket
Lemma~\ref{lem:blanket}, which gives
$P(\obs_{t+1:T}\given\state_t=i,\obs_{1:t})=P(\obs_{t+1:T}\given\state_t=i)$,
\[
  \fwd_t(i)\bwd_t(i)
  = P(\obs_{1:t},\state_t=i)\,P(\obs_{t+1:T}\given\state_t=i)
  = P(\obs_{1:t},\state_t=i)\,P(\obs_{t+1:T}\given\state_t=i,\obs_{1:t})
  = P(\obs_{1:T},\state_t=i).
\]
Dividing by $L=P(\obs_{1:T})$ gives the conditional \eqref{eq:gamma}.
\end{proof}

\begin{proposition}[Pairwise (transition) posterior]\label{prop:xi}
The joint posterior of consecutive states is, for $t=1,\dots,T-1$,
\begin{equation}\label{eq:xi}
  \xi_t(i,j)\eqdef P(\state_t=i,\state_{t+1}=j\given\obs_{1:T},\theta)
  = \frac{\fwd_t(i)\,A_{ij}\,b_j(\obs_{t+1})\,\bwd_{t+1}(j)}{L}.
\end{equation}
Consistency holds: $\sum_{j}\xi_t(i,j)=\gamma_t(i)$.
\end{proposition}

\begin{proof}
Analogous to Proposition~\ref{prop:gamma}, splitting the path at the
$t\!\to\!t\!+\!1$ transition and using Lemma~\ref{lem:blanket} to detach the
future $\obs_{t+1:T}$ from the past given $\state_{t+1}$:
$P(\obs_{1:T},\state_t=i,\state_{t+1}=j)
=\fwd_t(i)A_{ij}b_j(\obs_{t+1})\bwd_{t+1}(j)$, then normalise by $L$.
\end{proof}

\subsection{Numerical realisation: the log domain}\label{sec:logspace}

The forward variables are products of many densities and probabilities; they
decay to zero faster than double-precision floating point can represent
(underflow) within a few hundred steps. The implementation therefore carries
\emph{logarithms} of all quantities and replaces sums of probabilities by the
\emph{log-sum-exp} operation.

\begin{definition}[Log-sum-exp]
For $x=(x_1,\dots,x_n)\in\R^n$,
\begin{equation}\label{eq:lse}
  \operatorname{LSE}(x) \eqdef \log\sum_{k=1}^n e^{x_k}
  = m + \log\sum_{k=1}^n e^{x_k-m},
  \qquad m \eqdef \max_k x_k .
\end{equation}
\end{definition}
The right-hand identity is the numerically stable form actually evaluated: after
subtracting the maximum $m$, the largest exponent is $0$, so $e^{x_k-m}\in(0,1]$
and no term overflows; at least one term equals $1$, so the sum is bounded below
by $1$ and its logarithm never underflows (individual small terms may still round
to $0$, harmlessly --- it is the \emph{sum} that is protected). With $\ell\fwd_t(i)\eqdef\log\fwd_t(i)$ etc., the forward
recursion \eqref{eq:fwd-rec} becomes
\begin{equation}
  \ell\fwd_t(j) = \log b_j(\obs_t)
  + \operatorname{LSE}_i\!\big(\ell\fwd_{t-1}(i)+\log A_{ij}\big),
\end{equation}
and the Gaussian log-density is computed directly,
\begin{equation}
  \log b_k(\obs_t) = -\tfrac12\Big(\log(2\pi)+\log\sigma_k^2
                     + \frac{(\obs_t-\mu_k)^2}{\sigma_k^2}\Big),
\end{equation}
never as the logarithm of an exponentiated quantity. The smoothed posterior in
log form reads $\log\gamma_t(i)=\ell\fwd_t(i)+\ell\bwd_t(i)-\ell L$ with
$\ell L=\operatorname{LSE}_i \ell\fwd_T(i)$.

% ===========================================================================
\section{Learning: Baum--Welch as Expectation--Maximisation}\label{sec:em}
% ===========================================================================

We now estimate $\theta$ by (local) maximum likelihood,
$\hat\theta=\argmax_\theta \log L(\theta)$. Because the states are unobserved,
the log-likelihood $\log L(\theta)=\log\sum_{\state_{1:T}} P(\state_{1:T},\obs_{1:T}\given\theta)$
contains a logarithm of a sum and has no closed-form maximiser. The
Baum--Welch algorithm is the specialisation to HMMs of the
Expectation--Maximisation (EM) algorithm, which sidesteps this by iterating two
steps that are each tractable.

\subsection{The EM principle}\label{sec:em-principle}

Let $s\eqdef\state_{1:T}$ denote the latent path and $r\eqdef\obs_{1:T}$ the
data. For \emph{any} distribution $q$ over the latent paths, write the
\emph{evidence lower bound} (ELBO)
\begin{equation}\label{eq:elbo}
  \mathcal{F}(q,\theta) \eqdef \sum_{s} q(s)\,\log\frac{P(s,r\given\theta)}{q(s)} .
\end{equation}
A one-line manipulation gives the central identity
\begin{equation}\label{eq:em-identity}
  \log L(\theta) = \mathcal{F}(q,\theta) + \KL\!\big(q \,\big\|\, P(\cdot\given r,\theta)\big),
\end{equation}
where $\KL(q\|p)=\sum_s q(s)\log\frac{q(s)}{p(s)}\ge 0$ is the
Kullback--Leibler divergence, equal to zero iff $q=p$. Equation
\eqref{eq:em-identity} holds because
$\KL(q\|P(\cdot\given r,\theta))=\sum_s q(s)\log\frac{q(s)P(r\given\theta)}{P(s,r\given\theta)}
=\log L(\theta)-\mathcal F(q,\theta)$, using $P(s\given r,\theta)=P(s,r\given\theta)/P(r\given\theta)$.

EM exploits \eqref{eq:em-identity} by coordinate ascent on $\mathcal F$:
\begin{description}[leftmargin=2.2em,style=nextline]
  \item[E-step] At the current estimate $\theta^{(n)}$, maximise $\mathcal F$ over
        $q$. Since $\log L(\theta^{(n)})$ is fixed, \eqref{eq:em-identity} is
        maximised by driving the KL term to zero, i.e.\ by choosing
        \[
          q^{(n)}(s) = P(s\given r,\theta^{(n)}),
        \]
        the smoothed posterior over paths --- which the forward--backward
        algorithm gives us through its marginals $\gamma,\xi$.
  \item[M-step] Maximise $\mathcal F(q^{(n)},\theta)$ over $\theta$. Expanding the
        ELBO \eqref{eq:elbo} as a sum splits it into a $\theta$-dependent and a
        $\theta$-independent piece,
        \begin{equation}\label{eq:F-split}
          \mathcal F(q,\theta)
          = \underbrace{\sum_s q(s)\log P(s,r\given\theta)}_{\displaystyle Q(\theta\,\mid\,\theta^{(n)})}
          \;\underbrace{-\sum_s q(s)\log q(s)}_{\displaystyle H(q),\ \text{independent of }\theta},
        \end{equation}
        so maximising $\mathcal F(q^{(n)},\theta)$ over $\theta$ is equivalent to
        maximising the \emph{expected complete-data log-likelihood}
        \begin{equation}\label{eq:Qfun}
          Q(\theta\given\theta^{(n)}) \eqdef \E_{s\sim q^{(n)}}\!\big[\log P(s,r\given\theta)\big].
        \end{equation}
\end{description}

\begin{proposition}[Monotonicity]\label{prop:monotone}
The EM iterates never decrease the likelihood:
$\log L(\theta^{(n+1)})\ge \log L(\theta^{(n)})$.
\end{proposition}

\begin{proof}
With $q^{(n)}=P(\cdot\given r,\theta^{(n)})$ the KL term in
\eqref{eq:em-identity} vanishes, so
$\log L(\theta^{(n)})=\mathcal F(q^{(n)},\theta^{(n)})$. The M-step gives
$\mathcal F(q^{(n)},\theta^{(n+1)})\ge \mathcal F(q^{(n)},\theta^{(n)})$. Finally
\eqref{eq:em-identity} at $\theta^{(n+1)}$ with the \emph{same} $q^{(n)}$ reads
$\log L(\theta^{(n+1)})=\mathcal F(q^{(n)},\theta^{(n+1)})+\KL(\cdots)\ge
\mathcal F(q^{(n)},\theta^{(n+1)})\ge \log L(\theta^{(n)})$, the last KL being
$\ge 0$. Chaining the inequalities proves the claim.
\end{proof}

Proposition~\ref{prop:monotone} guarantees convergence of the likelihood
sequence (it is non-decreasing and, for our bounded-below variance model with a
floor, bounded above); it does \emph{not} guarantee convergence to the global
maximum. This is why we use multiple restarts (Section~\ref{sec:decisions}).

\subsection{The expected complete-data log-likelihood for an HMM}

Taking the logarithm of the factorised joint \eqref{eq:joint} turns products
into sums and exposes the parameters linearly through indicator functions:
\begin{equation}
  \log P(s,r\given\theta)
  = \sum_{i}\ind\{\state_1=i\}\log\pi_i
  + \sum_{t=2}^{T}\sum_{i,j}\ind\{\state_{t-1}=i,\state_t=j\}\log A_{ij}
  + \sum_{t=1}^{T}\sum_{i}\ind\{\state_t=i\}\log b_i(\obs_t).
\end{equation}
Taking the expectation under $q^{(n)}=P(\cdot\given r,\theta^{(n)})$ and using
$\E[\ind\{\state_t=i\}]=\gamma_t(i)$ and
$\E[\ind\{\state_{t-1}=i,\state_t=j\}]=\xi_{t-1}(i,j)$ gives a transition term
$\sum_{t=2}^{T}\xi_{t-1}(i,j)\log A_{ij}$; reindexing $t\mapsto t-1$ so the sum
runs over the $T-1$ transitions $t=1,\dots,T-1$ (the range on which $\xi_t$ is
defined, \eqref{eq:xi}) yields
\begin{equation}\label{eq:Q-hmm}
  Q(\theta\given\theta^{(n)})
  = \underbrace{\sum_{i}\gamma_1(i)\log\pi_i}_{\text{initial}}
  + \underbrace{\sum_{t=1}^{T-1}\sum_{i,j}\xi_t(i,j)\log A_{ij}}_{\text{transitions}}
  + \underbrace{\sum_{t=1}^{T}\sum_{i}\gamma_t(i)\log b_i(\obs_t)}_{\text{emissions}} .
\end{equation}
Crucially, the three groups of parameters $(\bm\pi)$, $(A)$, $(\mu,\sigma^2)$
appear in \emph{separate} additive terms, so the M-step decouples into three
independent maximisations.

\subsection{The M-step updates}

\paragraph{Initial distribution.}
Maximise $\sum_i\gamma_1(i)\log\pi_i$ subject to $\sum_i\pi_i=1$. With Lagrange
multiplier $\lambda$, the stationarity condition
$\partial_{\pi_i}\!\big[\sum_i\gamma_1(i)\log\pi_i+\lambda(1-\sum_i\pi_i)\big]=0$
gives $\gamma_1(i)/\pi_i=\lambda$, hence $\pi_i=\gamma_1(i)/\lambda$ with
$\lambda=\sum_k\gamma_1(k)$ fixed by the constraint. Since $\gamma_1(\cdot)$ is
itself a probability distribution over the states at $t=1$ (Proposition~\ref{prop:gamma}),
$\sum_k\gamma_1(k)=1$, so $\lambda=1$ and
\begin{equation}\label{eq:mstep-pi}
  \boxed{\;\pi_i^{(n+1)} = \gamma_1(i).\;}
\end{equation}
(The implementation divides by $\sum_k\gamma_1(k)$ explicitly rather than
relying on the exact identity, which is more robust to floating-point error;
the two agree because the sum is $1$.)

\paragraph{Transition matrix.}
For each row $i$, maximise $\sum_j\big(\sum_t\xi_t(i,j)\big)\log A_{ij}$ subject to
$\sum_j A_{ij}=1$. The identical Lagrange argument gives
$A_{ij}\propto\sum_t\xi_t(i,j)$, and using $\sum_j\xi_t(i,j)=\gamma_t(i)$ to
normalise,
\begin{equation}\label{eq:mstep-A}
  \boxed{\;A_{ij}^{(n+1)}
  = \frac{\sum_{t=1}^{T-1}\xi_t(i,j)}{\sum_{t=1}^{T-1}\gamma_t(i)}.\;}
\end{equation}
The numerator is the expected number of $i\!\to\!j$ transitions; the denominator
is the expected number of visits to $i$ from which a transition occurs --- the
estimator is exactly an ``expected count'' ratio.

\paragraph{Gaussian emissions.}
Only the emission term of \eqref{eq:Q-hmm} involves $(\mu_i,\sigma_i^2)$:
\[
  Q_i(\mu_i,\sigma_i^2)
  = \sum_{t=1}^{T}\gamma_t(i)\Big(-\tfrac12\log(2\pi)
    -\tfrac12\log\sigma_i^2 - \frac{(\obs_t-\mu_i)^2}{2\sigma_i^2}\Big).
\]
Setting $\partial Q_i/\partial\mu_i=0$:
$\sum_t\gamma_t(i)(\obs_t-\mu_i)/\sigma_i^2=0\Rightarrow$
\begin{equation}\label{eq:mstep-mu}
  \boxed{\;\mu_i^{(n+1)} = \frac{\sum_{t}\gamma_t(i)\,\obs_t}{\sum_{t}\gamma_t(i)}.\;}
\end{equation}
Setting $\partial Q_i/\partial\sigma_i^2=0$:
$\sum_t\gamma_t(i)\big(-\tfrac{1}{2\sigma_i^2}+\tfrac{(\obs_t-\mu_i)^2}{2\sigma_i^4}\big)=0\Rightarrow$
\begin{equation}\label{eq:mstep-sig}
  \boxed{\;(\sigma_i^2)^{(n+1)}
  = \frac{\sum_{t}\gamma_t(i)\,(\obs_t-\mu_i^{(n+1)})^2}{\sum_{t}\gamma_t(i)}.\;}
\end{equation}
These are exactly \emph{responsibility-weighted} sample mean and variance, with
the soft assignment $\gamma_t(i)$ playing the role of a fractional membership of
observation $t$ in state $i$. Writing $N_i\eqdef\sum_t\gamma_t(i)$ for the
expected occupancy of state $i$, the updates are
$\mu_i=N_i^{-1}\sum_t\gamma_t(i)\obs_t$ and
$\sigma_i^2=N_i^{-1}\sum_t\gamma_t(i)(\obs_t-\mu_i)^2$.

\subsection{The algorithm}

\begin{quote}\itshape
\textbf{Baum--Welch (one EM sweep).} Given $\theta^{(n)}$:
\begin{enumerate}[label=\arabic*.,leftmargin=2em]
  \item \textbf{E-step.} Run the forward (\ref{prop:fwd}) and backward
        (\ref{prop:bwd}) recursions in log-space; form $\gamma_t(i)$
        (\ref{prop:gamma}) and $\xi_t(i,j)$ (\ref{prop:xi}); record
        $\log L(\theta^{(n)})$.
  \item \textbf{M-step.} Update $\bm\pi,A,\mu,\sigma^2$ by
        \eqref{eq:mstep-pi}--\eqref{eq:mstep-sig}; apply the variance floor and
        any sign/zero constraints (Section~\ref{sec:decisions}).
  \item \textbf{Convergence.} Stop when
        $\log L(\theta^{(n+1)})-\log L(\theta^{(n)}) < \texttt{tol}$ or the
        iteration cap is reached.
\end{enumerate}
\end{quote}
One sweep costs $O(\nstate^2 T)$, dominated by the forward--backward pass.

% ===========================================================================
\section{Modelling decisions and their justification}\label{sec:decisions}
% ===========================================================================

The bare EM recursions above are textbook. What makes the estimator usable on
real return data is a set of deliberate choices, each addressing a concrete
failure mode.

\subsection{Why three Gaussian states}
Three states is the smallest number that can encode the qualitative trichotomy
(\emph{down / sideways / up}) that practitioners reason about, while keeping the
parameter count low ($\nstate^2-\nstate=6$ free transition probabilities,
$2\nstate=6$ emission parameters, $\nstate-1=2$ initial probabilities). A
Gaussian emission is the maximum-entropy choice given a state-specific mean and
variance, and it makes the M-step closed-form. Its known weakness --- thin tails
relative to real returns --- is discussed in Section~\ref{sec:extensions}.

\subsection{The variance floor: avoiding the likelihood singularity}\label{sec:floor}
The Gaussian-emission likelihood is \emph{unbounded above}. If a state's mean
$\mu_i$ coincides with a single observation $\obs_\tau$ and its variance
$\sigma_i^2\to 0$, then $b_i(\obs_\tau)=\big(2\pi\sigma_i^2\big)^{-1/2}\to\infty$.
The term of the path-sum \eqref{eq:like} in which $\state_\tau=i$ then contains
this diverging factor while all its remaining factors (the transitions and the
other emissions $b_{\state_u}(\obs_u)$, $u\ne\tau$) stay bounded below by a
positive constant as $\sigma_i^2\to0$; hence that single term, and with it
$L(\theta)$, diverges to $+\infty$. EM can march toward such a degenerate spike,
``explaining'' one day perfectly and collapsing the state. We forbid this by
flooring the variance after each M-step,
\begin{equation}
  \sigma_i^2 \leftarrow \max\!\big(\sigma_i^2,\ \varepsilon\big),
  \qquad \varepsilon = \texttt{var\_floor} \;(=10^{-8}\text{ by default}),
\end{equation}
which keeps every emission density bounded and the optimisation well-posed.

\subsection{Non-convexity and multiple restarts}
The log-likelihood surface is multimodal; EM converges only to a local maximum
(Proposition~\ref{prop:monotone}), and which one depends on the initial
$\theta^{(0)}$. We therefore run \texttt{n\_restarts} independent fits from
randomised initialisations and keep the one with the highest $\log L$. Means are
initialised by spreading them across empirical return quantiles (so distinct
restarts explore distinct mode assignments), variances around the sample
variance, and the transition matrix is initialised \emph{sticky}.

\subsection{Sticky transition initialisation}
Financial regimes persist for weeks or months, not single days. We initialise
$A$ with large diagonal mass,
\begin{equation}
  A^{(0)}_{ii}=0.9,\qquad A^{(0)}_{ij}=\frac{0.1}{\nstate-1}\ (j\ne i),
\end{equation}
encoding the prior belief that ``tomorrow's regime is probably today's''. This
both speeds convergence and steers EM away from implausible rapidly-switching
solutions.

\subsection{Identifiability: breaking label-switching by ordering}\label{sec:label}
The likelihood \eqref{eq:like} is \emph{invariant} under any permutation of the
state labels: relabelling states and permuting $(\bm\pi,A,\mu,\sigma^2)$
accordingly leaves $L$ unchanged. Hence there are $\nstate!$ equivalent global
modes, and ``state 1'' from one fit need not be the bear state. We remove this
ambiguity by adopting the canonical ordering
\begin{equation}
  \mu_1 \le \mu_2 \le \mu_3,
\end{equation}
i.e.\ after fitting we sort the states by mean ascending and permute all
parameters (and $A$ by the corresponding symmetric reindexing
$A\mapsto A[\sigma,\sigma]$). Index $1$ is then always \emph{bear} (lowest mean),
$2$ \emph{neutral}, $3$ \emph{bull}. This is what lets downstream code refer to
\texttt{P(bull)} unambiguously.

The ordering breaks the symmetry \emph{uniquely} precisely when the means are
distinct, $\mu_1<\mu_2<\mu_3$, which holds almost surely for continuous return
data. On the measure-zero event of an exact tie $\mu_i=\mu_j$ (with
$\sigma_i^2\ne\sigma_j^2$) the ordering by mean alone is not unique; the
implementation then breaks the residual tie by original index
(\texttt{np.argsort}), a deterministic but arbitrary convention. A fully
canonical rule would sort lexicographically on $(\mu_k,\sigma_k^2)$.

\subsection{Optional economic constraints}
Two optional constraints inject domain knowledge directly into the M-step:
\begin{itemize}[leftmargin=1.6em]
  \item \textbf{Pinned neutral mean} (\texttt{pin\_neutral\_zero}): force
        $\mu_{\text{neutral}}=0$ after \eqref{eq:mstep-mu}, encoding ``the
        sideways regime has no drift''. This is a hard equality constraint on
        one coordinate; the remaining updates are unchanged.
  \item \textbf{Sign constraint} (\texttt{enforce\_sign}): project the extreme
        means onto the economically admissible half-lines,
        $\mu_{\text{bear}}\leftarrow\min(\mu_{\text{bear}},0)$ and
        $\mu_{\text{bull}}\leftarrow\max(\mu_{\text{bull}},0)$, so the bear
        (bull) regime cannot end up with a positive (negative) drift.
\end{itemize}
Both act on the mean coordinate only, after the unconstrained update. The
monotonicity guarantee of Proposition~\ref{prop:monotone} is stated for the
\emph{unconstrained} M-step, so it does not transfer automatically under a
projection; what rescues it here is concavity. Each emission objective
$Q_i(\mu_i,\sigma_i^2)$ is, as a function of $\mu_i$, a negative quadratic and
hence strictly concave with unconstrained maximiser \eqref{eq:mstep-mu}. For a
strictly concave function on an interval, the maximiser over a closed feasible
subinterval is the point of that subinterval nearest the unconstrained
maximiser --- exactly what the clip $\min(\cdot,0)$ / $\max(\cdot,0)$ (or the
equality $\mu_{\text{neutral}}=0$) computes. The projected value is therefore the
\emph{constrained} maximiser of $Q$ over the feasible set, so the M-step still
maximises $Q$ subject to the constraint and the monotone-ascent guarantee
survives. (The sign clip is applied to the current lowest- and highest-mean
states, which under our ordering are the bear and bull states.)

% ===========================================================================
\section{Decoding: the Viterbi algorithm}\label{sec:viterbi}
% ===========================================================================

Whereas $\gamma_t$ gives the marginally most-likely state at each $t$
\emph{separately}, decoding asks for the single most-likely \emph{path}
\begin{equation}
  s^\star = \argmax_{s_{1:T}} P(\state_{1:T}=s_{1:T}\given\obs_{1:T},\theta)
          = \argmax_{s_{1:T}} P(\state_{1:T}=s_{1:T},\obs_{1:T}\given\theta),
\end{equation}
the second equality holding because the data term $P(\obs_{1:T})$ does not depend
on the path. The Viterbi algorithm solves this by the same dynamic-programming
factorisation, with $\max$ in place of $\sum$.

\begin{definition}[Viterbi variable]
\begin{equation}
  \delta_t(i) \eqdef \max_{s_{1:t-1}}
  P(\state_{1:t-1}=s_{1:t-1},\,\state_t=i,\,\obs_{1:t}\given\theta),
\end{equation}
the probability of the best path \emph{ending} in state $i$ at time $t$.
\end{definition}

\begin{proposition}[Viterbi recursion]
\begin{align}
  \delta_1(i) &= \pi_i\, b_i(\obs_1),\\
  \delta_t(j) &= \Big(\max_{i}\ \delta_{t-1}(i)\,A_{ij}\Big)\, b_j(\obs_t),
  \qquad
  \psi_t(j) = \argmax_{i}\ \delta_{t-1}(i)\,A_{ij},
\end{align}
with the optimal path recovered by back-tracking from
$s^\star_T=\argmax_i\delta_T(i)$ via $s^\star_t=\psi_{t+1}(s^\star_{t+1})$.
\end{proposition}

\begin{proof}
The recursion is Bellman's principle of optimality. By the joint factorisation
\eqref{eq:joint}, the value of a path ending in $\state_t=j$ factorises as
(value of its length-$(t-1)$ prefix ending in some $\state_{t-1}=i$)
$\times\,A_{ij}\,b_j(\obs_t)$, where the last two factors depend on the prefix
only through its endpoint $i$. Hence, \emph{exchange argument:} if an optimal
length-$t$ path ending in $j$ had a prefix that were not itself the best
length-$(t-1)$ path ending in its own endpoint $i$, we could substitute that
better prefix --- leaving $A_{ij}b_j(\obs_t)$ unchanged --- and strictly increase
the total, contradicting optimality. Therefore the optimal prefix is
$\delta_{t-1}(i)$ for the appropriate $i$, and maximising the product over the
predecessor $i$ gives the stated recursion; $\psi_t(j)$ stores the arg\,max so
the maximiser can be reconstructed.
\end{proof}

As with the forward pass, the computation is carried out on $\log\delta$ to avoid
underflow, turning products into sums and the $\max$ into a $\max$ of sums; this
is \emph{exact} (the logarithm is monotone, so it commutes with $\max$). Viterbi
is used for offline analysis --- e.g.\ painting the historical regime timeline ---
not for live trading, which must be causal.

% ===========================================================================
\section{Online inference: filtering and prediction}\label{sec:filtering}
% ===========================================================================

For trading we may use \emph{only past and present} data at each decision point;
using $\gamma_t$ (which conditions on the whole series, including the future)
would constitute look-ahead. The causal analogue of $\gamma_t$ is the
\emph{filtered belief}.

\begin{definition}[Filtered belief]
\begin{equation}
  \hat\gamma_t(i) \eqdef P(\state_t=i\given\obs_{1:t},\theta)
  = \frac{\fwd_t(i)}{\sum_{j}\fwd_t(j)} .
\end{equation}
\end{definition}
This is just the forward variable normalised across states at time $t$ --- it uses
no backward pass and hence no future information. In log-space the code computes
$\log\hat\gamma_t(i)=\ell\fwd_t(i)-\operatorname{LSE}_j\ell\fwd_t(j)$.

\paragraph{One-step state prediction.}
Given the belief at $t$, the predicted state distribution for $t+1$ propagates
through the transition matrix. Writing the belief as a column vector
$\hat\gamma_t\in\R^{\nstate}$,
\begin{equation}\label{eq:predict-state}
  p_{t+1}(j) \eqdef P(\state_{t+1}=j\given\obs_{1:t})
  = \sum_i \hat\gamma_t(i)\,A_{ij}
  = \big(A\Tr \hat\gamma_t\big)_j .
\end{equation}
(The transpose appears because $A$ is row-stochastic: $A_{ij}$ maps \emph{from}
$i$ \emph{to} $j$, so summing over the originating index $i$ is multiplication by
$A\Tr$.)

\paragraph{Predictive moments of tomorrow's return.}
Conditional on $\obs_{1:t}$, the next return $\obs_{t+1}$ is distributed as a
\emph{mixture of Gaussians} with weights $p_{t+1}(j)$:
\begin{equation}
  \obs_{t+1}\given\obs_{1:t} \ \sim\ \sum_{j=1}^{\nstate} p_{t+1}(j)\,\Normal(\mu_j,\sigma_j^2).
\end{equation}
Its first two predictive moments follow from the tower property. The mean is
\begin{equation}\label{eq:pred-mean}
  m_{t+1}\eqdef\E[\obs_{t+1}\given\obs_{1:t}] = \sum_j p_{t+1}(j)\,\mu_j ,
\end{equation}
and, using $\E[X^2]=\Var(X)+\E[X]^2$ within each component,
$\E[\obs_{t+1}^2\given\obs_{1:t}]=\sum_j p_{t+1}(j)\big(\sigma_j^2+\mu_j^2\big)$, so the
predictive variance is
\begin{equation}\label{eq:pred-var}
  v_{t+1}\eqdef\Var(\obs_{t+1}\given\obs_{1:t})
  = \sum_{j} p_{t+1}(j)\big(\sigma_j^2+\mu_j^2\big) \;-\; m_{t+1}^2 .
\end{equation}
Equation \eqref{eq:pred-var} is the \emph{law of total variance} with the regime
$\state_{t+1}$ as the conditioning variable, the relevant law being the one-step
predictive $\state_{t+1}\sim p_{t+1}$:
\begin{equation}
  \Var(\obs_{t+1}\given\obs_{1:t})
  = \E_{\state_{t+1}\sim p_{t+1}}\!\big[\,\sigma^2_{\state_{t+1}}\,\big]
  + \Var_{\state_{t+1}\sim p_{t+1}}\!\big(\mu_{\state_{t+1}}\big),
\end{equation}
i.e.\ the expected within-regime variance plus the across-regime dispersion of
means, both taken under $p_{t+1}$. (Both terms are conditional on $\obs_{1:t}$
through $p_{t+1}$; the inner $\E$ and $\Var$ average over the predicted regime.)
This is the single most important output for sizing: in a confident
low-volatility bull belief $v_{t+1}$ is small, whereas a belief that mixes a
high-variance bear state inflates $v_{t+1}$ even when the directional mean
$m_{t+1}$ is ambiguous.

% ===========================================================================
\section{From belief to position}\label{sec:trading}
% ===========================================================================

A trading rule is a map from the predictive summary $(p_{t+1},m_{t+1},v_{t+1})$
to a target position $w_{t+1}\in[-L,L]$, where $w=+1$ denotes fully long, $w=-1$
fully short, and $L$ is a leverage cap. Three rules are provided.

\subsection{Threshold rule}
The simplest rule acts only on directional confidence:
\begin{equation}
  w_{t+1} =
  \begin{cases}
    +1, & p_{t+1}(\text{bull}) > \tau,\\
    -1, & p_{t+1}(\text{bear}) > \tau,\\
    \phantom{+}0, & \text{otherwise,}
  \end{cases}
\end{equation}
for a confidence threshold $\tau\in(\tfrac12,1)$. It is robust and easy to reason
about, but ignores the magnitude of the expected edge and the predictive
variance.

\subsection{Mean--variance (fractional-Kelly) rule}\label{sec:kelly}
The rule that carries most of the model's value sizes the position by the
\emph{risk-adjusted} expected return. Consider a one-period mean--variance
objective for a position $w$ in an asset with predictive mean $m_{t+1}$ and
variance $v_{t+1}$,
\begin{equation}
  U(w) = w\,m_{t+1} - \tfrac{\lambda}{2}\, w^2 v_{t+1},
\end{equation}
where $\lambda>0$ is a risk-aversion coefficient. Setting $U'(w)=m_{t+1}-\lambda
w v_{t+1}=0$ gives the optimum
\begin{equation}\label{eq:kelly}
  \boxed{\;w_{t+1}^\star = \frac{m_{t+1}}{\lambda\, v_{t+1}},\;}
  \qquad\text{then clip to } [-L,L].
\end{equation}
Equation~\eqref{eq:kelly} is the \emph{exact} optimum of the stated
mean--variance utility $U$. It is closely related to --- but not identical with
--- the Kelly criterion, and it is worth being precise about the difference, as
the two are routinely conflated. The Kelly (growth-optimal) fraction maximises
the expected log-growth $\E[\log(1+w\,\obs_{t+1})]$. Two caveats apply for
Gaussian returns:
\begin{enumerate}[label=(\roman*),leftmargin=2em]
  \item \textbf{The objective is degenerate.} A Gaussian has full support on
        $\R$, so for any $w\ne 0$ the event $\{1+w\obs_{t+1}<0\}$ has positive
        probability and $\E[\log(1+w\obs_{t+1})]=-\infty$. Literal log-growth is
        therefore ill-defined here; the operative criterion is the
        \emph{quadratic} utility $U$, not log-utility.
  \item \textbf{$m/v$ is an approximation.} Expanding
        $\log(1+x)\approx x-\tfrac12 x^2$ gives
        $\E[\log(1+w\obs_{t+1})]\approx w m_{t+1}-\tfrac12 w^2(v_{t+1}+m_{t+1}^2)$,
        whose maximiser is $w^\star = m_{t+1}/(v_{t+1}+m_{t+1}^2)$. Only after the
        further approximation $m_{t+1}^2\ll v_{t+1}$ --- valid for daily returns,
        where the squared drift is tiny next to the variance --- does this reduce
        to $m_{t+1}/v_{t+1}$.
\end{enumerate}
With those qualifications, \eqref{eq:kelly} at $\lambda=1$ is the
small-return approximation to full Kelly, and $\lambda>1$ is \emph{fractional}
Kelly (a deliberately scaled-down bet). The economically important feature is
robust to all of this: the appearance of $v_{t+1}$ in the denominator is exactly
why the regime model is useful for risk control --- even with an unchanged
directional view $m_{t+1}$, the position automatically shrinks when the belief
admits a high-variance (bear/stressed) regime. The clip enforces the leverage
cap.

\subsection{Smooth (tanh) rule}
A continuous, bounded alternative maps the expected return through a hyperbolic
tangent,
\begin{equation}
  w_{t+1} = \tanh\!\Big(\frac{m_{t+1}}{c}\Big),
\end{equation}
where $c>0$ is the return scale at which the position reaches
$\tanh(1)\approx0.76$ of full size. It avoids the discontinuity of the threshold
rule and needs no variance estimate, at the cost of not penalising risk
directly.

\subsection{Transaction-cost control: the no-trade band}
Rebalancing every day to a slightly different target incurs transaction costs
that can erase a thin edge. We impose hysteresis: letting $w^{\text{prev}}$ be the
position currently held and $w^{\text{tgt}}$ the freshly computed target,
\begin{equation}
  w_{t+1} =
  \begin{cases}
    w^{\text{prev}}, & |w^{\text{tgt}}-w^{\text{prev}}| < \text{band},\\
    w^{\text{tgt}}, & \text{otherwise.}
  \end{cases}
\end{equation}
The band should be set so that a rebalance fires only when the change in target
position is large enough that the expected gain exceeds the round-trip cost of
trading it.

% ===========================================================================
\section{Out-of-sample evaluation: walk-forward backtesting}\label{sec:backtest}
% ===========================================================================

A backtest must never let information from the future leak into a past decision.
The harness enforces this by a \emph{rolling-window} (walk-forward) protocol.
Fix a training window length $W$ and a refit cadence $\rho$. For each trading day
$t>W$:
\begin{enumerate}[label=\arabic*.,leftmargin=2em]
  \item \textbf{Fit (periodically).} If $t-W$ is a multiple of $\rho$, refit the
        HMM by Baum--Welch on the trailing window $\obs_{t-W:t-1}$ only. (Daily
        refitting adds noise and cost; monthly, $\rho\approx 21$, is typical.)
  \item \textbf{Filter.} Compute the filtered belief $\hat\gamma_{t-1}$ on the
        same trailing history, then the predictive
        $(p_t,m_t,v_t)$ via \eqref{eq:predict-state}--\eqref{eq:pred-var}.
  \item \textbf{Size.} Apply the chosen rule (Section~\ref{sec:trading}) and the
        no-trade band to obtain the position $w_t$ \emph{held into} day $t$.
  \item \textbf{Realise P\&L.} The net strategy return on day $t$ is
        \begin{equation}
          \Pi_t = w_t\,\obs_t \;-\; \kappa\,|w_t-w_{t-1}|,
        \end{equation}
        where $\kappa$ is the per-unit-turnover transaction cost (in the same
        units as returns) and $|w_t-w_{t-1}|$ is the turnover.
\end{enumerate}
Because the parameters and the belief used to choose $w_t$ depend only on data up
to $t-1$, while the return $\obs_t$ is realised \emph{after} the position is set,
the protocol is causal by construction. The reported diagnostics --- annualised
return and volatility, Sharpe ratio, maximum drawdown, Calmar ratio, directional
hit rate, average gross exposure, annual turnover and cumulative cost --- are all
computed on the out-of-sample series $(\Pi_t)$.

\begin{remark}[An honesty check]
On structureless data, a \emph{correct} causal backtest should show no edge (a
Sharpe near zero or negative after costs). A backtest that reports a large Sharpe
on noise is almost always evidence of look-ahead leakage, not of skill. The
walk-forward construction above is the standard guard against that failure.
\end{remark}

% ===========================================================================
\section{Limitations and extensions}\label{sec:extensions}
% ===========================================================================

The model is a deliberately simple baseline. The main avenues for improvement,
roughly in order of expected impact:
\begin{itemize}[leftmargin=1.6em]
  \item \textbf{Heavy tails.} Replace the Gaussian emissions by Student-$t$
        emissions to capture the fat tails of daily returns; the M-step becomes
        an inner EM for the $t$ scale/degrees-of-freedom but the chain machinery
        is unchanged.
  \item \textbf{Conditional heteroskedasticity.} Let each state carry a GARCH
        recursion (a Markov-switching GARCH) so that volatility clusters
        \emph{within} a regime, not only across regimes.
  \item \textbf{Autoregressive emissions.} Allow $\E[\obs_t\given\state_t,\obs_{t-1}]$
        to depend on the previous return (Markov-switching AR), capturing
        short-horizon momentum or reversal inside a regime.
  \item \textbf{Duration modelling.} A plain HMM implies geometrically
        distributed regime durations; a hidden \emph{semi}-Markov model lets the
        dwell-time distribution be estimated.
  \item \textbf{Model selection.} Choose $\nstate$ and the emission family by an
        information criterion (BIC/AIC) on a validation window rather than fixing
        $\nstate=3$ a priori.
  \item \textbf{Multivariate / cross-asset states.} Drive a single regime chain
        from a vector of returns (TTF, HH, Brent jointly) with multivariate
        Gaussian emissions, so a regime is a market-wide rather than per-asset
        object.
\end{itemize}

% ===========================================================================
\section{Symbol dictionary and code map}\label{sec:symbols}
% ===========================================================================

Table~\ref{tab:symbols} maps the mathematical notation to the implementation in
\texttt{src/lng\_desk/hmm/}.

\begin{table}[h]
\centering
\small
\begin{tabular}{@{}lll@{}}
\toprule
Symbol & Meaning & Code \\
\midrule
$\nstate$ & number of hidden states ($=3$) & \texttt{n\_states} \\
$\state_t,\ \obs_t$ & hidden state, observed return at $t$ & latent / \texttt{r[t]} \\
$\bm\pi$ & initial state distribution & \texttt{HMMParams.pi} \\
$A$ & row-stochastic transition matrix & \texttt{HMMParams.A} \\
$\mu_k,\ \sigma_k^2$ & emission mean, variance of state $k$ & \texttt{HMMParams.mu}, \texttt{.sigma2} \\
$b_k(\obs)$ & Gaussian emission density & \texttt{\_gaussian\_logpdf} (log) \\
$\fwd_t(i),\ \bwd_t(i)$ & forward / backward variables & \texttt{\_forward\_backward} \\
$\gamma_t(i)$ & smoothed state posterior & \texttt{gamma} \\
$\xi_t(i,j)$ & pairwise transition posterior & \texttt{xi\_sum} (summed over $t$) \\
$\operatorname{LSE}$ & log-sum-exp & \texttt{\_logsumexp} \\
$\hat\gamma_t$ & filtered belief $P(\state_t\given\obs_{1:t})$ & \texttt{filter} \\
$p_{t+1}=A\Tr\hat\gamma_t$ & one-step state prediction & \texttt{predict\_next\_state} \\
$m_{t+1},\ v_{t+1}$ & predictive return mean / variance & \texttt{predict\_return\_moments} \\
$s^\star$ & Viterbi MAP path & \texttt{viterbi} \\
$\varepsilon$ & variance floor & \texttt{var\_floor} \\
$\lambda,\ L$ & risk aversion, leverage cap & \texttt{risk\_aversion}, \texttt{leverage\_cap} \\
$\tau,\ c$ & threshold, tanh scale & \texttt{tau}, \texttt{scale} \\
$W,\ \rho,\ \kappa$ & train window, refit cadence, cost & \texttt{train\_window}, \texttt{refit\_every}, \texttt{cost\_bps} \\
\bottomrule
\end{tabular}
\caption{Notation-to-code dictionary.}
\label{tab:symbols}
\end{table}

\paragraph{Summary.} The pipeline is: \emph{fit} parameters by multi-start
Baum--Welch on a trailing window (Section~\ref{sec:em}), with the variance floor
and label ordering of Section~\ref{sec:decisions}; \emph{filter} causally to get
the regime belief (Section~\ref{sec:filtering}); \emph{predict} tomorrow's return
mean and variance \eqref{eq:pred-mean}--\eqref{eq:pred-var}; \emph{size} the
position by a risk-aware rule \eqref{eq:kelly}; and \emph{evaluate} strictly out
of sample (Section~\ref{sec:backtest}). Every step reduces to the same
dynamic-programming factorisation of the Markov chain that makes the likelihood
\eqref{eq:like} computable in linear time.

\end{document}
